% Sketches de firmware TeleCon — qué .ino grabar.
% Build: pdflatex -output-directory=../build telecon_firmware_sketches.tex
% PDF entregado (asset de Android): ../build/firmware_sketches_esp.pdf
\documentclass[12pt,a4paper]{article}

\usepackage[margin=2.4cm]{geometry}
\usepackage[T1]{fontenc}
\usepackage[utf8]{inputenc}
\usepackage[spanish]{babel}
\usepackage{lmodern}
\usepackage{booktabs}
\usepackage{hyperref}
\usepackage{xcolor}
\usepackage{colortbl}
\usepackage{array}
\usepackage{enumitem}
\usepackage{needspace}
\usepackage{etoolbox}

\preto\section{\Needspace{8\baselineskip}}

\hypersetup{
  colorlinks=true,
  linkcolor=black,
  urlcolor=blue,
  pdftitle={Sketches de firmware TeleCon},
  pdfauthor={TeleCon}
}

\definecolor{TblHead}{HTML}{0F6F73}
\definecolor{TblChan}{HTML}{D5F5E3}
\definecolor{TblSw}{HTML}{FCF3CF}
\definecolor{TblTelem}{HTML}{D4E6F1}
\definecolor{TblKnob}{HTML}{F5CBA7}

\newcommand{\thc}[1]{\textcolor{white}{\bfseries #1}}
\newcommand{\chip}[2]{\colorbox{#1}{\strut\,\scriptsize #2\,}}
\newcommand{\ui}[1]{\textsf{\textbf{#1}}}

\title{Sketches de firmware TeleCon\\[0.3em]
\large Qué \texttt{.ino} grabar}
\author{ESP32 DevKit por Bluetooth, y ESP32-CAM para video}
\date{Versión 1.0 \quad 2026}

\begin{document}
\maketitle

\begin{abstract}
La aplicación Android descarga un ZIP por cada conexión desde el
repositorio de firmware. Cada ZIP es un sketch de Arduino, y el sketch es
el \texttt{.ino} del mismo nombre. Esta nota nombra cada sketch de control
por Bluetooth y el sketch de la ESP32-CAM. La cámara solo se conecta con
una DevKit por Bluetooth y una ESP32-CAM aparte para el video.
\end{abstract}

{\sffamily\scriptsize
\chip{TblChan}{DevKit + Bluetooth}
\chip{TblTelem}{CAM solo video + DevKit}}

\tableofcontents
\vspace{1em}

\section{La única forma de conectar la cámara}

Hoy la imagen no sale de la misma placa que las palancas. Hacen falta
dos placas, y no hay otro montaje de cámara en la aplicación.

\begin{center}
\renewcommand{\arraystretch}{1.28}
\setlength{\tabcolsep}{5pt}
\begin{tabular}{>{\raggedright\arraybackslash}p{3.6cm}
                >{\raggedright\arraybackslash}p{5.4cm}
                >{\raggedright\arraybackslash}p{5.2cm}}
\rowcolor{TblHead}
\thc{Papel} & \thc{Placa} & \thc{Sketch} \\
\rowcolor{TblChan}
Palancas & ESP32 DevKit & Uno de los tres \texttt{.ino} de Bluetooth \\
\rowcolor{TblTelem}
Imagen & ESP32-CAM & \texttt{TeleCon\_CAM\_\-SoftAP\_Video.ino} \\
\end{tabular}
\end{center}

Grabe el sketch de video en la CAM y un sketch Bluetooth en la DevKit.
En el teléfono, deje el modo Bluetooth apuntado a la DevKit y active la
cámara superpuesta. El teléfono se une al Wi-Fi
\texttt{TeleCon-RC-CAM-Starter}, contraseña \texttt{telecon1234}, solo
para la imagen. Las palancas siguen por Bluetooth. Una sola ESP32-CAM
no las lleva, y el control no viaja por el Wi-Fi de la cámara.

Sin cámara basta la DevKit con uno de los tres sketches Bluetooth. Esa
placa no lleva video.

\section{Control Bluetooth en una ESP32 DevKit}

Placa en la aplicación: \ui{DevKit}. Esta placa no lleva cámara. Abra la
carpeta del sketch en el IDE de Arduino (la carpeta que contiene el
\texttt{.ino}) y súbalo a un módulo ESP32 Dev.

\begin{center}
\renewcommand{\arraystretch}{1.28}
\setlength{\tabcolsep}{4pt}
\begin{tabular}{>{\raggedright\arraybackslash}p{4.3cm}
                >{\raggedright\arraybackslash}p{4.6cm}
                >{\raggedright\arraybackslash}p{5.3cm}}
\rowcolor{TblHead}
\thc{ZIP / \texttt{.ino}} & \thc{Ajuste en el teléfono} & \thc{En el aire} \\
\rowcolor{TblChan}
\texttt{TeleCon\_Classic\_Simple} & Classic + Simple & Nombre \texttt{ESP32-TC-RC-BT-Simple}, PIN \texttt{1234}. Saludo \texttt{proto,simple}. Líneas de texto \texttt{RC:CTRL}, \texttt{RC:BTN}, \texttt{RC:DATA}, \texttt{RC:PLOT}. \\
\rowcolor{TblChan}
\texttt{TeleCon\_Classic\_Binary} & Classic + Binary & Nombre \texttt{ESP32-TC-RC-BT-Binary}. Saludo \texttt{proto,binary}. Tramas \texttt{AA 55}, \texttt{BB 66}, \texttt{CC 11/22/33}. \\
\rowcolor{TblChan}
\texttt{TeleCon\_BLE\_Binary} & BLE Binary & Nombre \texttt{ESP32-TC-RC-BLE-Binary} (Nordic UART). Las mismas tramas binarias que Classic Binary. \\
\end{tabular}
\end{center}

\paragraph{\texttt{TeleCon\_Classic\_Simple.ino}}
Un solo archivo para la primera placa. Ponga el código de motores y servos
en \texttt{onControl()} y \texttt{onButton()}, donde el archivo dice
\texttt{YOUR CODE HERE}. El sketch ya incluye un servo de dirección en el
GPIO 13: la palanca \texttt{rx} $= 0$ es un pulso de $1500$\,\textmu s.
Los ocho bytes de telemetría (CH1--CH8 en el teléfono) se eligen en
\texttt{kTelemetryChannels[]} dentro de este mismo archivo. Con
\texttt{ENABLE\_PLOT} en 1 la placa envía \texttt{RC:PLOT} cada 50\,ms
aproximadamente.
La guía de código de este sketch es
\texttt{telecon\_classic\_simple\_code.tex} en esta documentación.

\paragraph{\texttt{TeleCon\_Classic\_Binary.ino}}
El \texttt{.ino} solo arranca el Bluetooth y los módulos del vehículo.
Los pines, el PWM y los ocho canales de telemetría están en los otros
archivos de esa carpeta. Edite \texttt{pins.h} para su controlador y
\texttt{TelemetryChannels.h} para lo que informan CH1--CH8. El ajuste del
teléfono debe ser Classic + Binary antes de conectar. El formato del cable
es el documento del protocolo de control, no las líneas de texto del
sketch Simple.

\paragraph{\texttt{TeleCon\_BLE\_Binary.ino}}
Los mismos módulos de vehículo y la misma tabla
\texttt{TelemetryChannels.h} que Classic Binary. El enlace es BLE en lugar
de Classic SPP. Elija BLE Binary en la aplicación. No lo empareje como un
dispositivo Classic.

Estos tres sketches son también la mitad de control del montaje de dos
placas que sigue. Grabe solo uno de ellos en la DevKit.

\section{ESP32-CAM, solo video}

Úselo cuando una DevKit ya ejecuta uno de los sketches Bluetooth y quiere
imagen desde una segunda placa.

\paragraph{\texttt{TeleCon\_CAM\_SoftAP\_Video.ino}}
ZIP \texttt{TeleCon\_CAM\_SoftAP\_Video.zip}. Grábelo en una ESP32-CAM
(AI-Thinker). Este sketch sirve video y no acepta palancas.

\begin{itemize}[leftmargin=1.4cm]
  \item SoftAP \texttt{TeleCon-RC-CAM-Starter}, contraseña
        \texttt{telecon1234}, dirección \texttt{192.168.4.1}.
  \item Imagen: \texttt{http://192.168.4.1/stream} (MJPEG) o
        \texttt{/capture} (un JPEG).
  \item Calidad del flujo: \texttt{http://192.168.4.1/camconfig}.
  \item No hay puerto de control. Esta placa no recibe \texttt{RC:CONNECT}.
\end{itemize}

En el teléfono, active la cámara superpuesta y deje la conexión Bluetooth
apuntada a la DevKit (Classic Simple, Classic Binary o BLE Binary). El
teléfono usa Wi-Fi para la CAM y Bluetooth para las palancas al mismo
tiempo. El control de las palancas no va en la ESP32-CAM.

\section{Qué contiene el ZIP}

La aplicación descarga \texttt{zips/TeleCon\_<nombre>.zip}. Descomprímalo
y abra la carpeta que contiene el \texttt{.ino} del mismo nombre. Arduino
compila cada \texttt{.cpp} y \texttt{.h} de esa carpeta, así que un sketch
binario es el \texttt{.ino} más esos módulos. No vuelva a copiar
\texttt{SteerCenter}: los sketches de serie no guardan un centro de
dirección. El recorte queda en el teléfono, y \texttt{rx} $= 0$ es
$1500$\,\textmu s.

\end{document}
